\documentclass[11pt,a4paper]{article}

% --- Packages ---
\usepackage[utf8]{inputenc}
\usepackage[margin=1in]{geometry}
\usepackage{titlesec}
\usepackage{hyperref}
\usepackage{graphicx}
\usepackage{booktabs} 
\usepackage{listings} 
\usepackage{xcolor}

% --- Python Code Styling ---
\definecolor{codegreen}{rgb}{0,0.6,0}
\definecolor{codegray}{rgb}{0.5,0.5,0.5}
\lstset{
    language=Python,
    basicstyle=\ttfamily\small,
    keywordstyle=\color{blue},
    commentstyle=\color{codegreen},
    stringstyle=\color{orange},
    breaklines=true,
    numbers=left,
    numberstyle=\tiny\color{codegray},
    frame=single
}

% --- Title Information ---
\title{JC1503: Object-Oriented Programming \\ \huge Group Project: Library Management System}
\author{CS Group 5}
\date{\today}

\begin{document}

\maketitle

% --- Requirement: Individual Contributions on First Page ---
\section*{Individual Contributions}
\begin{table}[h]
    \centering
    \begin{tabular}{lp{6cm}p{5cm}}
        \toprule
        \textbf{Member Name} & \textbf{Responsibilities} & \textbf{Technical Implementation} \\
        \midrule
        Zhan Zhijian & Architecture \& Integration & \texttt{LibSys} Controller, \texttt{Storage} (JSON), Git Mgmt \\
        Zou Wanyan & Business Model \& OOP Design & \texttt{Resource}, \texttt{Book}, \texttt{User} classes, Custom Exceptions \\
        Qiu Dairui & Linear Data Structures & Doubly Linked List, Stack (LIFO) for Undo \\
        Zhang Yuting & Search \& Queueing Structures & Hash Table (Chaining), Queue (FIFO) for waitlist \\
        Liao Yijing & Core Algorithms & Binary Search Tree (Recursive Insert/Search) \\
        Huang Yueshi & Category Tree \& QA & Multi-way Category Tree, Unit Testing \\
        \bottomrule
    \end{tabular}
    \caption{Summary of individual member responsibilities.}
\end{table}

\vfill
\begin{center}
    \textit{Note: Total report length must not exceed 10 pages.}
\end{center}
\newpage

% --- Section 1: Problem and System Design ---
\section{Problem and System Design}
\label{sec:problem_design}

\subsection{Application Domain}
Our chosen scenario is a Library Management System designed to address the challenges of managing library resources such as books and magazines, as well as tracking user borrowing records. In real-world libraries, librarians struggle with manually locating items, managing inventory, and tracking operations history. This system digitizes these processes, providing a centralized platform to seamlessly add, search, borrow, and return items, while preserving a persistent record of all transactions.

\subsection{System Architecture}
The system is architected as a complete application rather than a collection of disjointed scripts. It operates on a central Controller model (\texttt{LibSys} class in \texttt{main.py}), which orchestrates interactions between the user interface (CLI) and the underlying data models. 
Users interact with a menu-driven CLI that translates input into operations on our custom data structures. These structures (a \texttt{HashTable} for users and a \texttt{BST} for books) reside in memory during runtime. A dedicated \texttt{Storage} module handles data persistence, serializing the structures via their \texttt{to\_list()} and \texttt{to\_dict()} methods to a JSON file upon system exit, and deserializing it upon startup. 

% --- Section 2: Technical Design ---
\section{Technical Design}
\label{sec:technical_design}

\subsection{Object-Oriented Modelling}
The system heavily utilizes Object-Oriented Programming (OOP) principles:
\begin{itemize}
    \item \textbf{Inheritance:} We defined a base \texttt{Resource} class containing common attributes (\texttt{resource\_id}, \texttt{title}) and core methods (\texttt{borrow\_item}). \texttt{Book} and \texttt{Magazine} classes inherit from \texttt{Resource}, adding specific attributes like \texttt{author} and \texttt{issue\_number}.
    \item \textbf{Composition (Component Reuse):} A major architectural decision was to deeply integrate our custom data structures into our domain models. For example, the \texttt{User} class does not use Python's built-in \texttt{list}; instead, it is composed of our custom \texttt{DoublyLinkedList} to maintain the borrowing record. Similarly, the \texttt{Resource} base class utilizes our custom \texttt{Queue} to manage a real-time waitlist for out-of-stock items.
    \item \textbf{Encapsulation \& Polymorphism:} Data attributes are strictly modified via dedicated methods, safely raising custom exceptions (e.g., \texttt{OutOfStockError}). Polymorphism is elegantly demonstrated in data persistence: \texttt{Book} and \texttt{Magazine} override the \texttt{to\_dict()} method (utilizing \texttt{super().to\_dict()}) and implement polymorphic \texttt{@classmethod from\_dict(cls)} factories, allowing dynamic serialization.
\end{itemize}

\subsection{Data Structures and Algorithms}
We implemented all required custom data structures to manage the system's state:
\begin{itemize}
    \item \textbf{Stack (LIFO):} Used in \texttt{main.py} to implement a robust \texttt{undo\_last\_action()} feature, recording every system action (add, borrow, return) so they can be securely reversed.
    \item \textbf{Queue (FIFO):} Used inside the \texttt{Resource} class to manage waitlists. When an item is returned, the system automatically \texttt{dequeue()}s the next user and assigns the book.
    \item \textbf{Binary Search Tree (BST):} Utilized to store and manage books, using the book title string comparison as the branching condition. It includes recursive methods for \texttt{_add_node} and \texttt{_find_node}.
    \item \textbf{Hash Table:} A custom Hash Table (capacity=100) was developed to store \texttt{User} objects. It uses the user ID string as the key and implements chaining (Singly Linked Nodes) for collision resolution.
\end{itemize}

\subsection{Efficiency and Complexity}
Our algorithmic choices present distinct Big-$O$ behaviors:
\begin{itemize}
    \item \textbf{Hash Table (User Lookup):} The hash function calculates the ASCII sum of the key modulo 100. In the average case, this provides $O(1)$ time complexity. However, in the worst-case scenario (e.g., for anagram keys that produce the same ASCII sum), collisions are resolved via a linked chain, degrading lookup to $O(n)$, where $n$ is the chain length.
    \item \textbf{Binary Search Tree (Book Lookup):} The BST provides an average time complexity of $O(\log n)$ for recursive operations. Since our BST does not implement self-balancing, the worst-case complexity could degrade to $O(n)$ if books are inserted in strictly alphabetical order.
\end{itemize}

% --- Section 3: Development and Testing ---
\section{Development and Testing}
\label{sec:dev_testing}

\subsection{Iterative Development}
Initially, we focused on establishing the core OOP models and verifying their interactions. Next, we replaced built-in Python lists with our custom data structures. To fulfill the persistence requirement without hardcoding, we implemented a \texttt{generate\_initial\_data()} method in \texttt{storage.py}. On the very first run, it dynamically generates 30 users, 40 books, and 20 magazines, and establishes initial borrowing relationships, seamlessly saving them to \texttt{library\_data.json}.

\subsection{Problem Solving}
\subsubsection{Type Mismatch in BST Integration}
A critical bug occurred when integrating the BST with the borrowing logic. The search method returned the \texttt{Node} wrapper instead of the \texttt{Resource} object, causing an \texttt{AttributeError}. 

\begin{lstlisting}[caption={BST Search Fix: Accessing Value}]
# Incorrect Implementation: Returns the Node object
# Result: AttributeError: 'Node' has no attribute 'borrow_item'
result = bst.search("Book_Title_1")
result.borrow_item() 

# Corrected Implementation (Zhan Zhijian & Liao Yijing)
node = bst.search("Book_Title_1")
if node:
    node.value.borrow_item()
\end{lstlisting}

\subsubsection{JSON Serialization of Custom Objects}
When implementing data persistence, \texttt{json.dump()} failed with a \texttt{TypeError} because it cannot natively serialize custom classes. We solved this by implementing \texttt{to\_dict()} and \texttt{to\_list()} methods across all models and structures, effectively converting complex objects into primitive dictionaries before storage.

\subsubsection{Git Merge Conflict Resolution}
During integration, a merge conflict occurred in \texttt{main.py} between the "Menu Logic" and "Persistence" branches. Our team leader (Zhan Zhijian) organized a sync-up to manually inspect the Git diff, ensuring the JSON \texttt{load\_data()} executed \textit{before} the main menu loop initialized.

\subsection{Testing Strategy}
Our strategy included unit tests for individual structures (e.g., verifying Hash Table collision chaining) and system-level end-to-end testing via the CLI. We extensively tested the \texttt{Stack} by creating users, borrowing books, and sequentially calling the \texttt{undo} function to ensure the in-memory state correctly reverted.

% --- Section 4: Reflection ---
\section{Reflection}
\label{sec:reflection}

\subsection{Critical Assessment}
The system successfully fulfills the core requirements through a robust Object-Oriented design. The clear separation between data models (\texttt{Node}, \texttt{HashNode}) and manager classes (\texttt{BST}, \texttt{HashTable}) makes the codebase logical. 

However, our current design has limitations regarding scalability:
\begin{enumerate}
    \item \textbf{Hash Table Scaling:} Our ASCII sum hash function generates guaranteed collisions for anagrams. With a hardcoded capacity of 100 and no dynamic resizing logic, exceeding 100 users will cause collision chains to grow indefinitely, severely degrading $O(1)$ performance.
    \item \textbf{Unbalanced BST:} The lack of an auto-balancing mechanism means that importing a pre-sorted list of books during startup could result in a skewed, line-like tree, turning our $O(\log n)$ searches into $O(n)$ linear traversals.
\end{enumerate}

\subsection{Future Improvements}
If development continued, we would upgrade the Hash Table to use a more uniform algorithm (like DJB2) and implement dynamic array resizing. For the BST, we would refactor it into an AVL Tree to guarantee strict $O(\log n)$ height. Furthermore, we would transition from a CLI to a Graphical User Interface (GUI) to enhance user experience.

\subsection{Learning Outcomes}
Our group gained profound hands-on experience with software architecture. Writing data structures from scratch exposed us to the underlying mechanics of memory management and pointers (e.g., \texttt{next}, \texttt{left}, \texttt{right}). Collaborating on interdependent modules taught us the critical importance of clear interfaces—ensuring that a data structure returns the correct payload is vital for seamless team integration.

\end{document}